\section{Desire, reconciliation and the promised end}
The first interpretation explains the change needed in the person who hopes to see God. Augustine's growing capacity and Aquinas's distinction between preparation and consummation answer John's combination of present adoption, future likeness and self-purification. The second explains the identity and conduct of the people Christ gathers. Bede's symbolic fullness and baptismal cleansing answer the vision's question about the multitude. Chrysostom's account of the reconciling Son explains why peacemaking belongs to sonship and why fidelity may still bring opposition.

The difference affects how purity and peace are understood. In the first interpretation, a heart learning to seek God without divided desire approaches the promised vision. In the second, that heart belongs to persons undertaking reconciliation among their neighbors. Augustine's account of interior peace and Chrysostom's account of active peacemaking are compatible, but one cannot replace the other without losing an object of Christian charity.

Augustine and Aquinas understand the promised earth as a future inheritance. Chrysostom includes present possession alongside future gifts.

Both accounts require grace, willing action and perseverance. Their common future is communion with God in perfected charity. The prayers ask for help on the way and describe the saints' continued concern for those still seeking salvation. That communion lets the congregation rejoice in the saints while still asking to grow in the love whose completion it celebrates.
